\begin{thebibliography}{99}
\bibitem[1]{published-1} Miklós Ajtai. Σ11 -formulae on finite structures. Annals of Pure and Applied Logic, 24(1):1–48, 1983.
\bibitem[2]{published-2} Andris Ambainis, Leonard J. Schulman, Amnon Ta-Shma, Umesh V. Vazirani, and Avi Wigderson. The quantum communication complexity of sampling. SIAM J. on Computing, 32(6):1570–1585, 2003.
\bibitem[3]{published-3} Sanjeev Arora and Boaz Barak. Computational Complexity. Cambridge University Press, 2009. A modern approach.
\bibitem[4]{published-4} Paul Beame, Russell Impagliazzo, and Srikanth Srinivasan. Approximating ACˆ0 by small height decision trees and a deterministic algorithm for \#ACˆ0SAT. In Proceedings of the 27th Conference on Computational Complexity, CCC 2012, Porto, Portugal, June 26-29, 2012, pages 117–125. IEEE Computer Society, 2012. doi:10.1109/CCC.2012.40.
\bibitem[5]{published-5} Chris Beck, Russell Impagliazzo, and Shachar Lovett. Large deviation bounds for decision trees and sampling lower bounds for AC0-circuits. Electronic Colloquium on Computational Complexity (ECCC), 19:42, 2012.
\bibitem[6]{published-6} Chris Beck, Russell Impagliazzo, and Shachar Lovett. Large deviation bounds for decision trees and sampling lower bounds for AC0-circuits. In IEEE Symp. on Foundations of Computer Science (FOCS), pages 101–110, 2012.
\bibitem[7]{published-7} Avraham Ben-Aroya, Dean Doron, and Amnon Ta-Shma. Explicit two-source extractors for near-logarithmic min-entropy. Electronic Colloquium on Computational Complexity (ECCC), 23:88, 2016. URL: http://eccc.hpi-web.de/report/2016/088, arXiv:TR16-088.
\bibitem[8]{published-8} Itai Benjamini, Gil Cohen, and Igor Shinkar. Bi-lipschitz bijection between the boolean cube and the hamming ball. In IEEE Symp. on Foundations of Computer Science (FOCS), 2014.
\bibitem[9]{published-9} Eshan Chattopadhyay, Jesse Goodman, and David Zuckerman. The space complexity of sampling. Electron. Colloquium Comput. Complex., page 106, 2021. URL: https://eccc. weizmann.ac.il/report/2021/106, arXiv:TR21-106.
\bibitem[10]{published-10} Eshan Chattopadhyay and David Zuckerman. Explicit two-source extractors and resilient functions. In ACM Symp. on the Theory of Computing (STOC), pages 670–683, 2016.
\bibitem[11]{published-11} Gil Cohen. Making the most of advice: New correlation breakers and their applications. In IEEE Symp. on Foundations of Computer Science (FOCS), pages 188–196, 2016. doi: 10.1109/FOCS.2016.28.
\bibitem[12]{published-12} Gil Cohen and Leonard J. Schulman. Extractors for near logarithmic min-entropy. Electronic Colloquium on Computational Complexity (ECCC), 23:14, 2016.
\bibitem[13]{published-13} Thomas Cover and Joy Thomas. Elements of Information Theory (Wiley Series in Telecom- munications and Signal Processing). Wiley-Interscience, 2006.
\bibitem[14]{published-14} Imre Csiszar and Janos Korner. Information Theory: Coding Theorems for Discrete Memoryless Systems. Academic Press, Inc., 1982.
\bibitem[15]{published-15} Anindya De and Thomas Watson. Extractors and lower bounds for locally samplable sources. In Workshop on Randomization and Computation (RANDOM), 2011.
\bibitem[16]{published-16} P. Erdős and R. Rado. Intersection theorems for systems of sets. J. London Math. Soc., 35:85–90, 1960. doi:10.1112/jlms/s1-35.1.85.
\bibitem[17]{published-17} Merrick L. Furst, James B. Saxe, and Michael Sipser. Parity, circuits, and the polynomial-time hierarchy. Mathematical Systems Theory, 17(1):13–27, 1984.
\bibitem[18]{published-18} Anna Gál and Peter Bro Miltersen. The cell probe complexity of succinct data structures. Theoretical Computer Science, 379(3):405–417, 2007.
\bibitem[19]{published-19} Alexander Golynski. Cell probe lower bounds for succinct data structures. In 20th ACM-SIAM Symp. on Discrete Algorithms (SODA), pages 625–634, 2009.
\bibitem[20]{published-20} Mika Göös and Thomas Watson. A lower bound for sampling disjoint sets. ACM Trans. Comput. Theory, 12(3):20:1–20:13, 2020. doi:10.1145/3404858.
\bibitem[21]{published-21} Aryeh Grinberg, Ronen Shaltiel, and Emanuele Viola.               Indistinguishability by ad- aptive procedures with advice, and lower bounds on hardness amplification proofs. In IEEE Symp. on Foundations of Computer Science (FOCS), 2018.                    Available at https://www.ccs.neu.edu/home/viola/papers/adaptivemajority.pdf.
\bibitem[22]{published-22} Johan Håstad. Computational limitations of small-depth circuits. MIT Press, 1987.
\bibitem[23]{published-23} Johan Håstad. On the correlation of parity and small-depth circuits. SIAM J. on Computing, 43(5):1699–1708, 2014.
\bibitem[24]{published-24} Russell Impagliazzo, William Matthews, and Ramamohan Paturi. A satisfiability algorithm for AC0 . In ACM-SIAM Symp. on Discrete Algorithms (SODA), pages 961–972, 2012.
\bibitem[25]{published-25} Stasys Jukna. Boolean Function Complexity: Advances and Frontiers. Springer, 2012.
\bibitem[26]{published-26} Xin Li. Improved two-source extractors, and affine extractors for polylogarithmic entropy. In IEEE Symp. on Foundations of Computer Science (FOCS), 2016.
\bibitem[27]{published-27} Mingmou Liu and Huacheng Yu. Lower bound for succinct range minimum query. In Konstantin Makarychev, Yury Makarychev, Madhur Tulsiani, Gautam Kamath, and Julia Chuzhoy, editors, ACM Symp. on the Theory of Computing (STOC), pages 1402–1415. ACM, 2020. doi:10.1145/3357713.3384260.
\bibitem[28]{published-28} Shachar Lovett and Emanuele Viola. Bounded-depth circuits cannot sample good codes. Computational Complexity, 21(2):245–266, 2012.
\bibitem[29]{published-29} Peter Bro Miltersen. Cell probe complexity – A survey, 1999. Invited talk/paper at Advances in Data Structures (Pre-conference workshop of FSTTCS’99).
\bibitem[30]{published-30} Ryan O’Donnell. Analysis of Boolean Functions. Cambridge University Press, 2014.
\bibitem[31]{published-31} Mihai Pǎtraşcu. Succincter. In 49th IEEE Symp. on Foundations of Computer Science (FOCS). IEEE, 2008.
\bibitem[32]{published-32} Mihai Patrascu and Mikkel Thorup. Time-space trade-offs for predecessor search. In Jon M. Kleinberg, editor, ACM Symp. on the Theory of Computing (STOC), pages 232–240. ACM, 2006. doi:10.1145/1132516.1132551.
\bibitem[33]{published-33} Mihai Pǎtraşcu and Emanuele Viola. Cell-probe lower bounds for succinct partial sums. In 21th ACM-SIAM Symp. on Discrete Algorithms (SODA), pages 117–122, 2010.
\bibitem[34]{published-34} Alexander A. Razborov. Pseudorandom generators hard for k-DNF resolution and polynomial calculus resolution. Ann. of Math., 181(2):415–472, 2015. doi:10.4007/annals.2015.181.2.1.
\bibitem[35]{published-35} Nathan Segerlind, Sam Buss, and Russell Impagliazzo. A switching lemma for small restrictions and lower bounds for k-DNF resolution. SIAM J. on Computing, 33(5):1171–1200, 2004.
\bibitem[36]{published-36} Mikkel Thorup. Mihai patrascu: Obituary and open problems. Bulletin of the EATCS, 109:7–13, 2013. URL: http://albcom.lsi.upc.edu/ojs/index.php/beatcs/article/view/163/176.
\bibitem[37]{published-37} Emanuele Viola. The Complexity of Distributions, Fall 2018 talk at the Simons Institute. https://www.youtube.com/watch?v=O78b085HE3w.
\bibitem[38]{published-38} Emanuele Viola. Bit-probe lower bounds for succinct data structures. SIAM J. on Computing, 41(6):1593–1604, 2012.
\bibitem[39]{published-39} Emanuele Viola. The complexity of distributions. SIAM J. on Computing, 41(1):191–218, 2012.
\bibitem[40]{published-40} Emanuele Viola. Extractors for turing-machine sources. In Workshop on Randomization and Computation (RANDOM), 2012.
\bibitem[41]{published-41} Emanuele Viola. Extractors for circuit sources. SIAM J. on Computing, 43(2):355–972, 2014.
\bibitem[42]{published-42} Emanuele Viola. Lower bounds for data structures with space close to maximum imply circuit lower bounds. Theory of Computing, 15:1–9, 2019. URL: https://theoryofcomputing.org/ articles/v015a018/v015a018.pdf.
\bibitem[43]{published-43} Emanuele Viola. Sampling lower bounds: boolean average-case and permutations. SIAM J. on Computing, 49(1), 2020. Available at https://www.ccs.neu.edu/home/viola/papers/ sampling-lower-bounds.pdf.
\bibitem[44]{published-44} Emanuele Viola, 2022. https://emanueleviola.wordpress.com/2022/09/14/myth-creation -the-switching-lemma/.
\bibitem[45]{published-45} Andrew Yao. Separating the polynomial-time hierarchy by oracles. In 26th IEEE Symp. on Foundations of Computer Science (FOCS), pages 1–10, 1985.
\bibitem[46]{published-46} Andrew Chi-Chih Yao. Should tables be sorted? J. of the ACM, 28(3):615–628, 1981. doi:10.1145/322261.322274.
\bibitem[47]{published-47} Huacheng Yu. Optimal succinct rank data structure via approximate nonnegative tensor decomposition. In Moses Charikar and Edith Cohen, editors, ACM Symp. on the Theory of Computing (STOC), pages 955–966. ACM, 2019. doi:10.1145/3313276.3316352.
\end{thebibliography}
